% 2017_tianjin_liberal_q4_fig1。
% 节点清单：开始 / 输入 N / N能被3整除?菱形 / N=N/3(分式) / N=N-1 /
%   N≤3?菱形 / 输出 N / 结束。
% 有向边：开始→输入→J1→菱形1；菱形1是(下)→N=N/3→J2→菱形2；
%   菱形1否(右)→N=N-1 顶部；N=N-1 底部→(左转)→J2；
%   菱形2是(下)→输出 N→结束；菱形2否(左)→左侧大回环→J1。
% 首轮后修复：N=N/3 框及 J2/菱形2 下移（菱形1→N=N/3 直段 6.2mm→7.2mm，
%   是标签与框净距 0.974mm→约1.5mm）；其余拓扑不动。
\begin{tikzpicture}[x=1cm,y=1cm]
\tikzset{
  flownode/.style={draw, line width=0.6pt, align=center, inner sep=3pt,
    execute at begin node={\setlength{\parindent}{0pt}}},
  st/.style={flownode, shape=rectangle, rounded corners=1.5mm, minimum width=26mm, minimum height=10mm, inner xsep=3mm, inner ysep=1.5mm},
  io/.style={flownode, trapezium, trapezium left angle=68, trapezium right angle=112,
    minimum width=3.0cm, minimum height=1.0cm},
  pr/.style={flownode, minimum width=3.0cm, minimum height=1.0cm},
  dc/.style={flownode, diamond, aspect=2, minimum width=3.4cm, inner sep=2pt},
  fe/.style={draw, line width=0.6pt, -{Stealth[length=2.8mm,width=2.2mm]}},
  fl/.style={draw, line width=0.6pt}
}
\node[st] (st0) at (0,0) {开始};
\node[io] (ioN) at (0,-1.75) {输入 $N$};
\coordinate (J1) at (0,-2.85);
\node[dc,minimum width=5.6cm,aspect=3] (dc1) at (0,-4.45) {$N$能被$3$整除?};
\node[pr,minimum width=3.2cm,minimum height=1.7cm] (prD) at (0,-6.95) {$N=\dfrac{N}{3}$};
\coordinate (J2) at (0,-8.45);
\node[pr,minimum width=3.4cm] (prM) at (4.2,-6.9) {$N=N-1$};
\node[dc,minimum width=4.0cm,aspect=2.2] (dc2) at (0,-9.95) {$N\leqslant 3?$};
\node[io,minimum width=2.8cm] (ioN2) at (0,-12.0) {输出 $N$};
\node[st] (st1) at (0,-13.45) {结束};
\draw[fe] (st0.south) -- (ioN.north);
\draw[fe] (ioN.south) -- (J1) -- (dc1.north);

\draw[fe] (dc1.south) -- (prD.north);
\node[right=1.5mm,xshift=-0.4813mm,yshift=-0.1421mm] at (0,-5.74) {是};
\draw[fe] (dc1.east) -- (4.2,-4.45) -- (prM.north);
\node[above=1.5mm,xshift=0.1053mm,yshift=-0.2980mm] at (3.5,-4.45) {否};
\draw[fe] (prD.south) -- (J2) -- (dc2.north);

\draw[fe] (prM.south) -- (4.2,-8.45) -- (J2);
\draw[fe] (dc2.south) -- (ioN2.north);
\node[right=1.5mm,xshift=-0.4813mm,yshift=0.2316mm] at (0,-11.18) {是};
\draw[fe] (ioN2.south) -- (st1.north);
\draw[fe] (dc2.west) -- (-4.9,-9.95) -- (-4.9,-2.85) -- (J1);
\node[above=2mm,xshift=-0.1054mm,yshift=-0.7977mm] at (-2.5,-9.95) {否};
\end{tikzpicture}
